% Source printed pp. 41-46, entire Section 1.7.
\section{构造射影模空间}\label{src:1.7}
\begin{epigraph}
在我眼中, 空间 $M_{g,\nu}$(连同群 $\operatorname{SL}(2)$)是我在数学中遇到的最美, 最迷人的对象.
\epigraphauthor{Alexander Grothendieck~\cite[第 $129$ 页]{Gro86}}
\end{epigraph}
引入代数叠的一个动机, 是确保给定的模问题 $\mathcal M$ 可表, 并配备万有族. 虽然许多几何问题可以(也可以说应该)直接在模叠 $\mathcal M$ 上研究, 但作出一种取舍往往十分便利: 牺牲万有族的存在性, 有时就能构造一个更熟悉的几何空间, 理想情况下是射影簇. 这使我们可以利用射影几何更丰富的工具(如双有理几何, 相交理论, Hodge 理论等)研究模问题.

我们着重介绍构造射影模空间的两种方法: 
\begin{enumerate}[(1)]
\item 几何不变量理论(GIT); 以及
\item 粗模空间／良模空间的内在构造.
\end{enumerate}
内在方法与外在方法之间有着优美的相互作用. GIT 的思想启发了下述内在方法六个步骤中的每一步; 反过来, 内在方法又加深了对 GIT 的理解. GIT 也与 $19$ 世纪的不变量理论紧密交织, 确定 GIT 半稳定轨迹本身就是一个有趣而重要的问题. 同时牢记这两种方法很有价值.

\subsection{GIT 方法}\label{src:1.7.1-subsection}
\begin{quote}
\begin{center}\textbf{GIT 策略概述}\end{center}
\begin{enumerate}[(A)]
\item\label{step:git-a} 将模叠 $\mathcal M$ 表示为子叠
\[
\mathcal M\subseteq[X/G],
\]
其中 $G$ 为约化群, $X\hookrightarrow\mathbb P(V)$ 为到 $G$-表示 $V$ 的射影化中的 $G$-等变嵌入.
\item\label{step:git-b} 证明点 $x\in X$ 为 GIT 半稳定, 当且仅当 $x\in\mathcal M$; 换言之, $\mathcal M=[X^{\mathrm{ss}}/G]$.
\end{enumerate}
\end{quote}
对步骤~\ref{step:git-a}, 通常可以通过参数化附加数据, 自然地将模问题刚性化. 对光滑曲线 $C$, 选取 $\Gamma(C,\Omega_C^{\otimes3})$ 的一组基, 就定义三重典范嵌入 $C\hookrightarrow\mathbb P^{5g-6}$(例~\bookref{1.6.13}); 或者, 对任意 $k\geq3$, 选取 $\Gamma(C,\Omega_C^{\otimes k})$ 的一组基, 就定义第 $k$ 重典范嵌入 $C\hookrightarrow\mathbb P^{(2k-1)(g-1)-1}$. 对固定光滑曲线 $C$ 上的向量丛 $E$, 选择充分大的整数 $m$ 后, 选取 $\Gamma(C,E(m))\cong\mathbb C^{P(m)}$ 的一组基, 就可将 $E$ 视为商 $\mathcal O_C(-m)^{P(m)}\twoheadrightarrow E$.

刚性化后的模问题应有一个可由射影簇 $X$ 表出的紧化: 在这两个例子中, 分别是 $\operatorname{Hilb}^P(\mathbb P^{5g-6})$ 和 $\operatorname{Quot}^P(\mathcal O_C(-m)^{P(m)})$. 附加数据的选择应由群 $G$ 的作用控制. 要使 GIT 方法成功, 需要 $G$ 约化, 且 $\mathcal M$ 为 $[X/G]$ 的子叠. 最后, 需要选择 $G$-等变嵌入 $X\hookrightarrow\mathbb P(V)$, 其中 $V$ 为有限维 $G$-表示; 等价地, 选择 $X$ 上某个丰沛线丛的 $G$-线性化.

步骤~\ref{step:git-b}~最困难: 必须证明 $\mathcal M$ 恰为 $[X/G]$ 中 GIT 半稳定点组成的开子叠. 利用 Hilbert--Mumford 判别法(\bookref{8.4.4}), 可将问题转化为: 点 $x\in X$ 表示模问题 $\mathcal M$ 的一个对象, 当且仅当对每个单参数子群 $\lambda:\mathbb G_m\to G$, Hilbert--Mumford 指标 $\mu(x,\lambda)\geq0$. 这往往将目标归约为一个可处理(但常常仍十分棘手)的组合问题.

GIT 商 $M:=X^{\mathrm{ss}}/\!/G$ 必为射影的. GIT 的一个优美特点是, 即使模叠 $\mathcal M$ 不紧, GIT 策略也能给出紧化!若 $\mathcal M$ 只有有限自同构, 等价地, 没有严格半稳定点, 则 $X^{\mathrm{ss}}\to M$ 为几何商, $\mathcal M=[X^{\mathrm{ss}}/G]\to M$ 为粗模空间. 若存在无穷自同构, 则 $X^{\mathrm{ss}}\to M$ 为良商, $\mathcal M\to M$ 为良模空间.

GIT 方法在第~\bookref{8}~章详细介绍. 我们将在 \S\bookref{6.8} 概述 $\overline M_g$ 的 GIT 构造, 并在 \S\bookref{9.6} 给出 $\cBun_{r,d}(C)$ 的完整 GIT 构造.

\subsection{内在方法}\label{src:1.7.2-subsection}
\begin{quote}
\begin{center}\textbf{通向射影模空间的六个步骤}\end{center}
\begin{enumerate}[(1)]
\item\label{step:intrinsic-1}\textbf{代数性: }将模叠 $\mathcal M$ 表示为一个更大的模叠 $\mathcal X$ 的子叠
\[
\mathcal M\subseteq\mathcal X.
\]
定义对象 $x\in\mathcal X$ 为\textit{半稳定的}, 当且仅当它属于 $\mathcal M$; 由此可以将 $\mathcal M$ 视为半稳定轨迹 $\mathcal X^{\mathrm{ss}}$. 证明 $\mathcal X$ 为 $\mathbb C$ 上局部有限型的代数叠.
\item\label{step:intrinsic-2}\textbf{半稳定性的开性: }证明半稳定性为开条件, 即 $\mathcal M=\mathcal X^{\mathrm{ss}}\subseteq\mathcal X$ 为开子叠.
\item\label{step:intrinsic-3}\textbf{半稳定性的有界性: }证明半稳定性有界, 即 $\mathcal M=\mathcal X^{\mathrm{ss}}$ 在 $\mathbb C$ 上为有限型.
\item\label{step:intrinsic-4}\textbf{半稳定约化: }证明 $\mathcal M$ 满足固有性的赋值判别法的存在性部分.
\item\label{step:intrinsic-5}\textbf{模空间的存在性: }证明存在细模空间／粗模空间／良模空间 $\mathcal M\to M$, 其中 $M$ 为固有代数空间.\footnote{花体 $\mathcal M$ 表示叠, 而正体 $M$ 表示空间. 全书都遵循这一约定.}
\item\label{step:intrinsic-6}\textbf{射影性: }证明 $\mathcal M$ 上的某个重言线丛下降为 $M$ 上的丰沛线丛, 即 $M$ 为射影的.
\end{enumerate}
\end{quote}
GIT 神奇地一次解决了所有这些步骤!在 GIT 方法的步骤~\ref{step:git-a}~中, 将模叠 $\mathcal M$ 表示为 $[X/G]$ 的子叠, 已经蕴含``有界性''. 由于 GIT 半稳定性总是开条件, 步骤~\ref{step:git-b}~中将 $\mathcal M$ 等同于半稳定轨迹 $[X^{\mathrm{ss}}/G]$, 就给出了``半稳定性的开性'', 从而给出 $\mathcal M$ 的``代数性''. 更令人惊叹的是, GIT 还蕴含其余各步: ``半稳定约化''``模空间的存在性''以及``射影性''.

\begin{enumerate}[步骤 1: ]
\item\textbf{代数性.} 许多模叠都有自然的扩大. 光滑曲线的叠 $\mathcal M_g$ 与稳定曲线的叠 $\overline{\mathcal M}_g$, 都包含于所有曲线的叠或极化曲线的叠. 光滑曲线上半稳定向量丛的叠, 包含于所有向量丛的叠, 或更大的所有凝聚层的叠. 通常较容易先证明扩大后的 $\mathcal X$ 为局部有限型代数叠, 再用步骤~\ref{step:intrinsic-2}~推得 $\mathcal M$ 本身也是代数叠.

首先需要定义叠 $\mathcal M$ 及其扩大 $\mathcal X$, 这意味着要指定对象族以及拉回和等同. 要检验 $\mathcal X$ 的代数性, 需要找出由概形给出的光滑覆盖 $U\to\mathcal X$. 许多情况下, 甚至可以证明 $\mathcal X$ 等同于商叠 $[U/G]$, 此时 $U\to[U/G]$ 就给出了一个呈示. 或者, 常常可以使用 Artin 判别准则(定理~\bookref{C.7.4})建立代数性; 这本质上归结为检验模问题的局部性质, 特别是需要理解形变与障碍理论.

\item\textbf{半稳定性的开性.} 这转化为如下条件: 对概形 $S$ 上的每个 $\mathcal X$ 的对象族 $\mathcal E$, 子集
\begin{equation}\label{src:1.7.1}
\{s\in S\mid\mathcal E_s\text{ 半稳定, 即 }\mathcal E_s\in\mathcal M\},
\end{equation}
是 $S$ 的开子集, 其中 $\mathcal E_s$ 为 $\mathcal E$ 沿 $\operatorname{Spec}\kappa(s)\to S$ 的拉回. 这正是包含映射 $\mathcal M=\mathcal X^{\mathrm{ss}}\hookrightarrow\mathcal X$ 可由开浸入表出的含义: 对每个映射 $S\to\mathcal X$(对应于族 $\mathcal E$), 纤维积 $\mathcal M\times_{\mathcal X}S$(在集合意义下等同于~\eqref{src:1.7.1})是 $S$ 的开子概形. 这一步保证 $\mathcal M$ 也是 $\mathbb C$ 上局部有限型的代数叠.

\item\textbf{半稳定性的有界性.} 称 $\mathbb C$ 上的代数叠 $\mathcal M$ \newterm{有界}{bounded}, 是指它为有限型. 由于步骤~\ref{step:intrinsic-2}~意味着 $\mathcal M$ 在 $\mathbb C$ 上局部有限型, 有界性就转化为 $\mathcal M$ 的拟紧性. 更具体地, 有界性等价于: 存在 $\mathbb C$ 上的有限型概形 $Z$ 以及 $Z$ 上的对象族 $\mathcal E$, 使 $\mathcal M$ 的每个对象 $E$ 都同构于某个(未必唯一的)$z\in Z$ 对应的 $\mathcal E_z$.

例如, $\mathcal M_g$ 有界, 但所有亏格为 $g$ 的固有曲线的叠, 以及所有光滑曲线(任意亏格)的叠 $\coprod_g\mathcal M_g$ 都无界. 对向量丛, 我们将证明固定秩和次数的半稳定向量丛的叠 $\cBun_{r,d}^{\mathrm{ss}}(C)$ 有界. \textit{reference} 固定秩和次数的所有向量丛的叠 $\cBun_{r,d}(C)$ 无界, 任意秩和次数的半稳定向量丛的叠也无界.

\item\textbf{半稳定约化.} 固有性的赋值判别法的存在性部分断言: 对每个离散赋值环 $R$(可以将它看成光滑曲线的局部模型), 其分式域为 $K$(或去掉一点的曲线), $K$ 上的每个对象 $\mathcal E^\times$, 在可能将 $R$ 换成离散赋值环的扩张后, 都延拓为 $R$ 上的对象族 $\mathcal E$. 换言之, 每个图
\begin{equation}\label{src:1.7.2}
\begin{tikzcd}
\operatorname{Spec}K\arrow[r,"\mathcal E^\times"]\arrow[d] & \mathcal M\\
\operatorname{Spec}R\arrow[ur,dashed,"\mathcal E"'] &
\end{tikzcd}
\end{equation}
在将 $R$ 换成扩张后, 都有一个延拓. 若 $R$ 上的延拓 $\mathcal E$ 还唯一, 则称 $\mathcal M$ 满足固有性的赋值判别法; 这蕴含固有性(定理~\bookref{4.8.7}), 特别地, 也蕴含分离性. 可以说, 赋值判别法在代数几何中的用处, 在模理论中体现得最为充分.

光滑曲线的模叠不紧, 也不满足赋值判别法的存在性部分.
\begin{figure}[H]
\centering\begin{tikzpicture}[x=1cm,y=1cm,>=Stealth]
\draw (-3,.85)--(3,.85);
\draw (-3,-.85)--(3,-.85);
\foreach \a in {-2.7,0,2.7} {
\draw plot[domain=-.85:.85,samples=70] ({\a+.14*sin(360*\x)},\x);
}
\foreach \a in {-1.5,1.5} {
\draw[blue,thick,domain=-1.65:1.65,samples=140]
plot ({\a+.28*(\x*\x-1)},{.28*\x*(\x*\x-1)});
}
\draw[->,thick] (0,-1.12)--(0,-1.82);
\draw[<->] (-3,-2.15)--(3,-2.15) node[right] {$\lambda$};
\foreach \a/\t in {-1.5/0,1.5/1} {
\draw (\a,-2.15)--(\a,-2.22) node[below] {$\t$};
}
\end{tikzpicture}

\caption{椭圆曲线族 $y^2z=x(x-z)(x-\lambda z)$ 在 $\lambda=0,1$ 处退化为结点三次曲线.}\label{src:1.7.3}
\end{figure}
射影簇当然是紧的, 并满足赋值判别法. 若希望构造射影模空间, 模叠就应满足赋值判别法的存在性部分. $\overline{\mathcal M}_g$ 的固有性由 Deligne 和 Mumford 在其影响深远的论文~\cite{DM69}~中首次证明. 我们将在 \S\bookref{6.5} 证明特征 $0$ 中的半稳定约化.

对于向量丛的模, 半稳定约化首先由 Mumford 和 Seshadri 作为 GIT 构造的推论证明~\cite{Ses67}. 后来 Langton 给出了内在的几何论证~\cite{Lan75}. 注意, 与稳定曲线不同, 叠 $\cBun_{r,d}^{\mathrm{ss}}(C)$ 不分离, 因为 $C_K$ 上的向量丛可能有几个不同构的 $C_R$ 上的延拓. 尽管如此, 模叠 $\cBun_{r,d}^{\mathrm{ss}}(C)$ 满足一种较弱的分离性, 称为 \newterm{$S$-完备性}{S-completeness}.

\item\textbf{模空间的存在性.} 我们希望构造一个代数空间, 使它尽可能好地逼近模叠. 这一步取决于模问题的自同构: 
\begin{itemize}
\item \textit{没有自同构: }此时模叠 $\mathcal M$ 已经是一个代数空间 $M$. 换言之, $M$ 为细模空间: 对概形 $S$, $S$ 上的对象与映射 $S\to M$ 之间有自然双射.
\item \textit{有限自同构: }必须证明 $\mathcal M$ 分离, 换言之, 除赋值判别法的存在性部分外, $\mathcal M$ 还满足唯一性部分. Keel--Mori 定理(定理~\bookref{5.4.6})便给出粗模空间 $\mathcal M\to M$ 的存在性, 其中 $M$ 为固有代数空间. 映射 $\mathcal M\to M$ 在 $\mathbb C$-点上诱导双射, 并满足如下万有性质: 任意其他到代数空间 $Y$ 的映射 $\mathcal M\to Y$, 都唯一地经 $M$ 分解.
\item \textit{约化自同构: }除证明 $\mathcal M$ 满足固有性的赋值判别法的存在性部分外, 还必须证明 $\mathcal M$ 满足另外两个赋值判别条件, 称为 \newterm{$\Theta$-完备性}{Theta-completeness}和 $S$-完备性. 它们要求: 在某些带有 $\mathbb G_m$-作用的去掉一点的曲面上, $\mathbb G_m$-等变对象族具有延拓; 见 \S\bookref{7.9.2-subsection}. 具有这些性质后, 良模空间存在性定理(\bookref{7.10.1})给出良模空间 $\mathcal M\to M$, 其中 $M$ 为固有代数空间.

映射 $\mathcal M\to M$ 不再在 $\mathbb C$-点上给出双射, 因为它等同闭包相交的点, 其方式与 GIT 中轨道闭包等价关系类似. 但 $\mathcal M\to M$ 确实在 $\mathcal M$ 的闭 $\mathbb C$-点(有时称为多重稳定对象)与 $M$ 的 $\mathbb C$-点之间诱导双射, 而且仍满足到代数空间的映射的万有性质.
\end{itemize}

\item\textbf{射影性.} 这通常是最困难的一步. 它要求深入理解模问题的几何, 有时还依赖双有理几何中的精细技术. Kollár 在~\cite{Kol90}~中引入一种检验簇的模叠的射影性的策略, 并用它检验 $\overline M_g$ 的射影性. 我们将在 \S\bookref{6.9} 介绍 Kollár 的方法. Faltings 在~\cite{Fal93}~中不使用 GIT 理论就构造了向量丛的射影模空间, 我们在 \S\bookref{9.6} 的构造中借鉴了他的若干思想.
\end{enumerate}
